% 编译建议：latexmk -xelatex main.tex

\usepackage[a4paper,margin=2.5cm,headheight=15pt]{geometry}
\usepackage{amsmath,amssymb,amsfonts,mathtools,bm}
\usepackage{booktabs,array,longtable,multirow}
\usepackage{enumitem}
\usepackage{graphicx}
\usepackage{xcolor}
\usepackage{caption}
\usepackage{subcaption}
\usepackage{float}
\usepackage{tikz}
\usetikzlibrary{arrows.meta,positioning,calc}
\usepackage[most]{tcolorbox}
\usepackage{algorithm2e}
\usepackage{listings}
\usepackage[
  backend=biber,
  style=gb7714-2015,
  sorting=nyt
]{biblatex}
\usepackage{hyperref}
\usepackage[nameinlink,noabbrev]{cleveref}

\addbibresource{references.bib}
\graphicspath{{figures/}}

\hypersetup{
  colorlinks=true,
  hypertexnames=false,
  linkcolor=blue!60!black,
  citecolor=green!45!black,
  urlcolor=purple!60!black,
  pdftitle={强化学习入门},
  pdfauthor={PatFang}
}

\setlist[itemize]{leftmargin=2em,itemsep=2pt,topsep=4pt}
\setlist[enumerate]{leftmargin=2em,itemsep=2pt,topsep=4pt}
\captionsetup{font=small,labelfont=bf,labelsep=quad}

\definecolor{noteBlue}{RGB}{35,95,170}
\definecolor{noteGreen}{RGB}{36,130,90}
\definecolor{noteOrange}{RGB}{190,105,30}
\definecolor{notePurple}{RGB}{105,70,155}
\definecolor{codeBg}{RGB}{248,248,248}
\definecolor{codeFrame}{RGB}{220,220,220}
\definecolor{codeKeyword}{RGB}{0,92,175}
\definecolor{codeComment}{RGB}{88,135,85}
\definecolor{codeString}{RGB}{160,70,70}

\tcbset{
  enhanced,
  breakable,
  arc=2mm,
  boxrule=0.6pt,
  left=1.5mm,
  right=1.5mm,
  top=1mm,
  bottom=1mm,
  colback=white,
  before upper={\setlength{\parindent}{2em}}
}

\newtcolorbox{definitionbox}[1][]{
  title={定义 - #1},
  colframe=noteBlue,
  colback=noteBlue!4
}

\newtcolorbox{keybox}[1][]{
  title={关键结论 - #1},
  colframe=noteGreen,
  colback=noteGreen!4
}

\newtcolorbox{warningbox}[1][]{
  title={注意 - #1},
  colframe=noteOrange,
  colback=noteOrange!6
}

\newtcolorbox{questionbox}[1][]{
  title={待思考 - #1},
  colframe=notePurple,
  colback=notePurple!5
}

\newenvironment{itemparagraph}{%
  \par\smallskip
  \begingroup
  \leftskip=\dimexpr\leftskip-\labelwidth-\labelsep\relax
  \setlength{\parindent}{2em}%
  \indent\ignorespaces
}{%
  \par\endgroup\smallskip
}

\SetKwInput{KwInput}{输入}
\SetKwInput{KwOutput}{输出}
\SetKw{KwInit}{初始化}
\SetKw{KwReturn}{返回}
\SetAlgoLined
\DontPrintSemicolon

\lstdefinestyle{pythonstyle}{
  language=Python,
  basicstyle=\ttfamily\small,
  keywordstyle=\color{codeKeyword}\bfseries,
  commentstyle=\color{codeComment},
  stringstyle=\color{codeString},
  backgroundcolor=\color{codeBg},
  frame=single,
  rulecolor=\color{codeFrame},
  breaklines=true,
  showstringspaces=false,
  tabsize=4,
  numbers=left,
  numberstyle=\tiny\color{gray}
}
\lstset{style=pythonstyle}

% 常用集合与算子
\newcommand{\R}{\mathbb{R}}
\newcommand{\E}{\mathbb{E}}
\newcommand{\Prob}{\mathbb{P}}
\newcommand{\1}{\mathbf{1}}
\newcommand{\KL}{\mathrm{KL}}
\newcommand{\Entropy}{\mathcal{H}}
\newcommand{\argmax}{\operatorname*{arg\,max}}
\newcommand{\argmin}{\operatorname*{arg\,min}}

% 强化学习符号
\newcommand{\states}{\mathcal{S}}
\newcommand{\actions}{\mathcal{A}}
\newcommand{\transition}{P}
\newcommand{\reward}{r}
\newcommand{\discount}{\gamma}
\newcommand{\policy}{\pi}
\newcommand{\trajectory}{\tau}
\newcommand{\replay}{\mathcal{D}}
\newcommand{\params}{\theta}
\newcommand{\oldparams}{\theta_{\mathrm{old}}}
\newcommand{\valuefn}{V^{\policy}}
\newcommand{\qvalue}{Q^{\policy}}
\newcommand{\advantage}{A^{\policy}}
\newcommand{\tdtarget}{y^{\mathrm{TD}}}
\newcommand{\tdelta}{\delta}

\newcommand{\figref}[1]{\cref{#1}}
\newcommand{\tabref}[1]{\cref{#1}}
\newcommand{\eqnref}[1]{\cref{#1}}
